\section{High-Fidelity Simulation of Rayleigh--Taylor Instability}

Direct Numerical Simulations are used as reference data because they resolve the relevant spatial and temporal scales of the flow without relying on turbulence closure models. In the context of RTI, they make it possible to observe the transition from early interface perturbations to non-linear bubble, spike, and mixing structures. Their main limitation is computational cost, which restricts the number of available realizations and motivates data-driven surrogate modeling.

DNS has become a central tool for turbulence research because it gives direct
access to the full resolved flow field and makes it possible to analyze
physical mechanisms without introducing an explicit turbulence model. Moin and
Mahesh describe DNS as a reference approach for studying turbulent flows, while
also emphasizing its strong dependence on computational resources and spatial
resolution \citep{MoinMahesh1998}. This limitation is particularly important in
the present work: RTI simulations are rich enough to describe multi-scale
mixing, but expensive enough that only a limited dataset can be generated.
Pope's reference text on turbulent flows also provides the physical background
for interpreting coherent structures, mixing, spectra, and statistical
quantities in turbulent regimes \citep{pope2000turbulent}.

